\section{Proceso de referencia y comparación entrópica}\label{sec:referencia}

El perfil continuo proporciona los pesos para construir colocaciones paso a paso.
El objetivo de esta sección es comparar la ley de las colocaciones así construidas,
condicionada al éxito, con la ley uniforme de la capa. La cota de entropía relativa
permitirá trasladar a esta última las estimaciones obtenidas durante la construcción.

Fijamos $K$ y $0\le h\le K$, y ponemos
\[
 q=n-h,\qquad N=2n-1,\qquad H=N-q=n+h-1 .
\]
En esta sección y en las dos siguientes, para $h\ge1$ el perfil se evalúa en el parámetro
objetivo $x_*=h/n$, fijo durante todo el proceso, y omitimos el subíndice $x_*$. Donde aparezcan
a la vez este perfil objetivo y el perfil lleno, como en la Sección~\ref{subsec:trayectoria},
escribiremos de nuevo el subíndice $x_*$ en el primero. La capa $h=0$
se trata por separado en la Sección~\ref{subsec:h0}.

\subsection{Hipergrafo aumentado y pesos}

Se añaden $n-1$ vértices \emph{ficticios} a las filas y otros tantos a las columnas, de modo que
cada una de las cuatro familias tiene $N$ vértices. Las \emph{aristas reales} son las $n^2$
casillas, con sus cuatro líneas físicas. Las \emph{aristas virtuales} son todas las cuaternas
con un vértice de cada familia, físico o ficticio. El color (real o virtual) forma parte de la
arista: una arista virtual es distinta de una real aunque tengan los mismos vértices.

En la Figura~\ref{fig:cruces-vacios}, las líneas de $(3,2)$ son $R_3$, $C_2$,
$D_1$ y $A_5$, pues $3-2=1$ y $3+2=5$.
Una elección real en esa casilla coloca una reina. Una elección virtual con
esas mismas cuatro líneas las desactiva en el proceso sin colocar ninguna
reina en el tablero. Esta distinción se usará al contar las reparaciones. Llamaremos \emph{piezas} a las
aristas de una
trayectoria, es decir, a las que elige el proceso, reales o virtuales según su color.

Para $h\ge1$, los pesos de vértice son $u_R(\bar r)=x_*r_R(\bar r)$ en las filas físicas,
$u_C(\bar c)=x_*r_C(\bar c)$ en las columnas físicas, $u_D$ y $u_A$ en las coordenadas inducidas
de diagonales y antidiagonales, y $1$ en los vértices ficticios. Todos son positivos. El peso de
una arista, real o virtual, es el producto $\varpi(e)=\prod_{v\in e}u_v$; por el
Teorema~\ref{thm:perfil}(a), toda arista real tiene peso entre $cx_*^2$ y $Cx_*^2$.

\subsection{El proceso}

\begin{definicion}[Proceso de referencia]\label{def:referencia}
Se preasigna, sin revelarla, una palabra uniforme con $q$ colores reales y $H$ virtuales. En
cada paso $t=1,\dots,N$ se revela sólo el siguiente color, se elige una arista disponible (con
sus cuatro vértices todavía activos) de ese color, con probabilidad proporcional a su peso, y se
desactivan sus cuatro vértices. Una pieza virtual no coloca ninguna reina. Si el normalizador
del color pedido, es decir, la suma de los pesos de sus aristas disponibles, es cero, el proceso pasa a un
estado de fallo. $\QQ_h$ es la ley de las
trayectorias y $S$ el evento de éxito (los $N$ pasos se completan).
\end{definicion}

En $S$, las $q$ aristas reales elegidas son casillas sin líneas comunes, es decir, una
colocación de $\Om_h(n)$: la \emph{colocación final}. Sea $\nu'_h$ su ley bajo
$\QQ_h(\cdot\mid S)$. En las elecciones reales puede cancelarse el factor común $x_*^2$ sin
cambiar la ley; no se cancela en las virtuales. La referencia no necesita tener éxito con
probabilidad próxima a uno.

\subsection{La ley de comparación}

\begin{definicion}[Ley de comparación]\label{def:comparacion}
Se elige $P\in\Om_h(n)$ uniforme; después, una completación virtual uniforme, es decir, un
emparejamiento perfecto por aristas virtuales de los $H$ vértices libres de cada familia (hay
$(H!)^3$); finalmente, un orden uniforme de las $N$ piezas. $\PP_h$ es la ley de la sucesión de
aristas resultante.
\end{definicion}

Representamos el orden uniforme mediante tiempos independientes y uniformes en $[0,1]$
asociados a las piezas; las piezas se revelan en orden creciente de tiempo. (La uniformidad del
tiempo de una pieza aislada no bastaría: se usa la independencia.) Estos tiempos independientes,
junto con la desigualdad de Jensen del Lema~\ref{lem:normalizadores}(iii), son el recurso de las
cotas por entropía de Radhakrishnan~\cite{Radhakrishnan} y de Linial y Luria~\cite{LinialLuria}.
La ley $\PP_h$ proyecta
exactamente en $\mu_h$, y toda trayectoria de su soporte tiene probabilidad positiva bajo
$\QQ_h$: la siguiente pieza del emparejamiento completo sigue disponible y tiene peso positivo.
En particular $\PP_h(S)=1$.

\subsection{La identidad entrópica}

Sean $W_R(t)$ y $W_V(t)$ los pesos totales de las aristas disponibles de cada color antes del
paso $t$, y $\col(t)$ el color revelado en ese paso.

\begin{lema}\label{lem:identidad}
Para $1\le h\le K$,
\[
 \KL{\PP_h}{\QQ_h}=\E_{\PP_h}\sum_{t=1}^N\log W_{\col(t)}(t)-\sum_v\log u_v-\log N_h-\log q!
 -4\log H!,
\]
donde $\sum_v$ recorre los $4N$ vértices.
\end{lema}

\begin{proof}
Sea $\omega$ una trayectoria del soporte de $\PP_h$. Entonces $\PP_h(\omega)=[N_h(H!)^3N!]^{-1}$.
Bajo $\QQ_h$, la palabra de colores tiene probabilidad $\binom Nq^{-1}$ y cada elección tiene
probabilidad $\varpi(e_t)/W_{\col(t)}(t)$. Como el emparejamiento completo cubre cada vértice
exactamente una vez, el producto de los pesos elegidos es $\prod_vu_v$, independiente del
emparejamiento y del orden. Por tanto
\[
 \log\frac{\PP_h(\omega)}{\QQ_h(\omega)}=\sum_t\log W_{\col(t)}(t)-\sum_v\log u_v
 -\log\bigl[N_h(H!)^3N!\bigr]+\log\binom Nq ,
\]
y $-\log[(H!)^3N!]+\log\binom Nq=-\log q!-4\log H!$. Se toma la esperanza bajo $\PP_h$.
\end{proof}
